\documentclass[11pt]{article}
\usepackage[T1]{fontenc}
\usepackage{lmodern,booktabs,amsmath,geometry,hyperref}
\geometry{margin=1in}
\hypersetup{hidelinks}
\title{When does a biodiversity recording peak survive missing dates and dataset exclusion?\\[3pt]\large A five-species audit of Balearica}
\author{Lluis Eriksson}
\date{20 November 2026}
\begin{document}
\maketitle
\begin{abstract}
For five named species, 12,251 Balearica occurrence rows generate 24,502 references to GBIF. Each row links to one occurrence and its dataset. Using distinct occurrences, we calculate which months can lead the recording counts if missing months are assigned arbitrarily. For the rock pigeon, April necessarily leads with all datasets included, yet removing one dataset changes the dated leader to May. For the seagrass \emph{Posidonia oceanica}, only August and October can lead. Checking all 54 undated references at GBIF resolves no month. One dataset contributes 45 of these records and no dated records. Its exclusion makes October certain. These conclusions concern the composition of public records and the information in their structured fields.
\end{abstract}
\section{Question and material}
A portal can show one distribution by source and another by month. Comparing their totals assumes that both count the same unit. Planning around a monthly maximum requires two further choices: how to treat incomplete dates and which contributors to retain. We ask how the answer changes under these choices in a single snapshot of Balearica, a biodiversity platform of the Universitat de les Illes Balears~\cite{bal,method}.

We chose three widespread species before inspecting monthly totals: a marine plant, a terrestrial tree and a bird. After finding that the source and monthly totals differed for all three, we added a common olive and pigeon before querying their counts. This purposive design examines specific counterexamples. It cannot estimate how often the pattern occurs across the Balearic biota. We queried the public occurrence list, its count and the source statistics by accepted species ID on 20 November 2026. We verified list completeness against the count and checked distinct internal and external occurrence IDs within each species. We used one GBIF dataset key per row for the exclusion analysis.

\section{Criteria}
Let $c_j$ be the number of dated records in month $j$ and $u$ the number of records without a month. Under unrestricted completion of the missing months, month $j$ can share first place if and only if
\[
 c_j+u\ \geq\ \max_k c_k .
\]
Necessity follows from the best possible increase for $j$. Sufficiency follows by placing all $u$ records in $j$. Month $j$ is guaranteed strictly first exactly when $c_j>\max_{k\ne j} c_k+u$: allocating all missing records to the strongest rival attains the worst case. An empty retained panel has no winner. Subtracting a dataset's twelve monthly counts and its undated count gives the same tests for its exclusion. A dataset whose records are all undated can therefore affect certainty while leaving every observed bar height unchanged.

\section{Findings}
All 12,251 rows contained two links with GBIF as provider: an occurrence key and a dataset key. In each species the number of links reproduces the source aggregate, while distinct occurrence--provider pairs reproduce the occurrence count. This describes the observed responses; the portal's database query was not inspected.
\begin{center}\small
\begin{tabular}{lrrrrl}
\toprule
Species & Rows & Datasets & Undated & Dated leader & Possible leaders\\
\midrule
\emph{T.\ merula} & 3 & 1 & 0 & tie & 1, 2, 4\\
\emph{P.\ oceanica} & 590 & 22 & 54 & 10 & 8, 10\\
\emph{Q.\ ilex} & 1,559 & 26 & 212 & 1 & all except 7\\
\emph{O.\ europaea} & 3,152 & 21 & 562 & 1 & all\\
\emph{C.\ livia} & 6,947 & 15 & 231 & 4 & 4\\
\bottomrule
\end{tabular}
\end{center}
Months run from January (1) to December (12). Dataset counts include contributors with no dated rows. The blackbird, \emph{Turdus merula}, has three records in three different months from the same dataset. Two of its three external occurrence links resolved to matching GBIF records. The other returned HTTP 404 at the time of the check.

The rock pigeon, \emph{Columba livia}, has 1,144 dated April records and 842 in May, the runner-up. Even assigning all 231 undated records to May leaves April ahead by 71. Removing dataset \texttt{4fa7b334-ce0d-4e88-aaae-2e0c138d049e} changes the leader among dated records to May. In the holm oak, \emph{Quercus ilex}, and the olive, \emph{Olea europaea}, removing the same dataset from each panel changes the dated leader from January to October. The accompanying tables identify every dataset and permit all single exclusions to be reproduced.

\section{Tracing the seagrass dates}
In the seagrass panel October leads August by 114 to 69 among dated records. Allocating 46 undated records to August and the remaining eight to January gives August a strict lead; allocating them all to October gives October a strict lead. To check whether these possibilities relied on information lost in transfer, we queried every undated GBIF occurrence key. All 54 returned the matching species and dataset. None supplied a numeric month. The structured event date was absent in 53 and year-only in one.

The Anthos dataset~\cite{anthos} contributes 45 of these 54 records and no dated rows. Its public dataset coverage gives a range of years. With that dataset omitted, all 536 dated records remain. Only nine months then require completion, leaving October at least $114-(69+9)=36$ ahead of every rival. This is a sensitivity calculation; the analysis does not justify discarding the Anthos records. Excluding a different dataset can instead change the \emph{dated} leader to August. More detailed primary collecting notes could further constrain the possible months.

\section{Interpretation, reproducibility and review context}
The pigeon and the seagrass illustrate different decisions. A recording peak may survive every assignment of incomplete dates and still change when a dataset is omitted. Conversely, a source without a single dated record can determine whether the peak is identifiable. The analysis provides no correction for variation in observation effort, habitat, years or contributor coverage. It measures recording patterns, with no inference about biological abundance or an optimal survey month.

The reproducibility package contains month-by-dataset counts, hashes and source URLs for the public occurrence lists, code for exact possible and guaranteed leaders, and an audit of the external structured date fields. The records may change after the snapshot. The counting inequalities are elementary. The contribution is the verified application and the identified records driving each conclusion.

\paragraph{Model assessments and review context.} Codex assisted with data extraction, analysis and drafting. Scientific assessments, if attached to a released version, apply to the exact PDF identified by its hash. This text records no independent assessment or peer-review decision.
\begin{thebibliography}{9}
\bibitem{bal} T. Cancellario and colleagues. Balearica: integrating taxonomic, distributional, ecological and genetic biodiversity data for the Balearic Islands. \href{https://doi.org/10.3897/BDJ.14.e209972}{doi:10.3897/BDJ.14.e209972}.
\bibitem{method} Centre Balear de Biodiversitat, Universitat de les Illes Balears. \href{https://balearica.uib.cat/method/getting_started/distribution.html}{Taxon distribution methods} and \href{https://balearica.uib.cat/api/docs}{Balearica API}.
\bibitem{anthos} CSIC--Real Jard\'in Bot\'anico. Anthos, Sistema de Informaci\'on de las Plantas de Espa\~na. \href{https://doi.org/10.15468/4wnutv}{doi:10.15468/4wnutv}.
\end{thebibliography}
\end{document}
